% El beneficio del café como diagrama de flujo — las tres estructuras.
%
% Diagrama propio en TikZ (sin librerías extra: los rombos y los óvalos se
% dibujan con `shape` básicos y coordenadas explícitas). Se tiñe con el color
% de la línea (milcGradeDark = verde lima de Pensamiento computacional).
%
% Muestra el proceso documentado por Cenicafé (Cartilla cafetera, cap. 20):
% despulpar el mismo día, fermentar, lavar, secar; y en medio, la prueba que
% decide el paso siguiente —«si la muestra suena a cascajo»—, con su lazo de
% repetición mientras todavía no suena.
%
% Las etiquetas de dentro enseñan (nombran la estructura de cada figura); la
% síntesis va en el pie de figura vía \guiaDiagrama, no aquí.
\begin{tikzpicture}[
    font=\sffamily,>=stealth,
    accion/.style={fill=milcGradeDark,text=white,rounded corners=3pt,
                   inner sep=4pt,align=center,font=\scriptsize\bfseries,text width=2.6cm},
    extremo/.style={draw=milcNegro,line width=0.7pt,rounded corners=9pt,
                    inner sep=4pt,align=center,font=\scriptsize\bfseries,text width=1.5cm},
    nota/.style={font=\scriptsize\itshape,milcNegro!70,align=left,text width=3.2cm}
  ]

  % ── Columna principal del flujo ──────────────────────────────────────
  \node[extremo] (ini) at (0,0) {Inicio};
  \node[accion]  (des) at (0,-1.25) {despulpar\\el mismo día};
  \node[accion]  (tan) at (0,-2.65) {poner en\\el tanque};

  % ── El rombo: la decisión ────────────────────────────────────────────
  \node (rombo) at (0,-4.5) {};
  \draw[fill=milcMostaza,draw=milcNegro,line width=0.6pt]
    (0,-3.55) -- (2.25,-4.5) -- (0,-5.45) -- (-2.25,-4.5) -- cycle;
  \node[align=center,font=\scriptsize\bfseries,milcNegro,text width=3.2cm] at (0,-4.5)
    {¿la muestra suena\\a cascajo?};

  \node[accion] (lav) at (0,-6.5) {lavar};
  \node[accion] (sec) at (0,-7.9) {secar};
  \node[extremo] (fin) at (0,-9.2) {Fin};

  % ── La rama del «no»: esperar y volver a probar ──────────────────────
  \node[accion] (esp) at (4.6,-4.5) {esperar\\una hora};

  % ── Flechas ──────────────────────────────────────────────────────────
  \draw[->,milcGradeDark,line width=1pt] (ini) -- (des);
  \draw[->,milcGradeDark,line width=1pt] (des) -- (tan);
  \draw[->,milcGradeDark,line width=1pt] (tan) -- (0,-3.55);
  \draw[->,milcGradeDark,line width=1pt] (0,-5.45) -- (lav);
  \draw[->,milcGradeDark,line width=1pt] (lav) -- (sec);
  \draw[->,milcGradeDark,line width=1pt] (sec) -- (fin);

  % salida «no» → esperar → regreso al rombo (el lazo)
  \draw[->,milcGradeDark,line width=1pt] (2.25,-4.5) -- (esp);
  \draw[->,milcGradeDark,line width=1pt]
    (esp) -- ++(0,1.9) -- (0,-2.6) -- (0,-3.55);

  \node[font=\scriptsize\bfseries,milcNegro] at (2.72,-4.26) {no};
  \node[font=\scriptsize\bfseries,milcNegro] at (0.42,-5.85) {sí};

  % ── Leyenda: qué estructura es cada figura ───────────────────────────
  \node[nota] at (-4.3,-1.25) {rectángulo:\\\textbf{secuencia}};
  \node[nota] at (-4.3,-4.5) {rombo:\\\textbf{decisión}};
  \node[nota] at (-4.3,-7.2) {flecha de regreso:\\\textbf{repetición}};

  \draw[milcNegro!35,line width=0.5pt,dashed] (-2.6,-1.25) -- (-2.1,-1.25);
  \draw[milcNegro!35,line width=0.5pt,dashed] (-2.6,-4.5) -- (-2.35,-4.5);
  \draw[milcNegro!35,line width=0.5pt,dashed] (-2.6,-7.2) -- (-1.4,-7.2);

\end{tikzpicture}
